% Dirichlet Round 05 English translation TeX
% Source: G. Lejeune Dirichlet's Werke, Band II (1897), Paper VI, title p. 49 and body pp. 51--66.
% Prepared 2026-06-02. Editable TeX. Translation checked against the Band II scan.
\documentclass[11pt]{article}
\usepackage[a4paper,margin=1in]{geometry}
\usepackage[T1]{fontenc}
\usepackage[utf8]{inputenc}
\usepackage{lmodern}
\usepackage{amsmath,amssymb}
\usepackage{microtype}
\usepackage{hyperref}
\hypersetup{colorlinks=true,linkcolor=black,urlcolor=black,citecolor=black}
\setlength{\parindent}{1.2em}
\setlength{\parskip}{0.2em}
\newcommand{\floor}[1]{\left[#1\right]}
\begin{document}

\begin{center}
{\Large G. Lejeune Dirichlet's Werke. II.}\par\vspace{1.2em}
{\Large VI.}\par\vspace{0.5em}
{\LARGE \textbf{On}}\par
{\LARGE \textbf{the Determination of Mean Values}}\par
{\LARGE \textbf{in Number Theory}}\par\vspace{1em}
{\large G. Lejeune Dirichlet}\par\vspace{0.5em}
{\small Proceedings of the Royal Prussian Academy, 1849, pp. 69--83.}
\end{center}

\begin{center}
\textbf{On the determination of mean values in number theory.}\par
{\small [Read in the Academy of Sciences on 9 August 1849\footnotemark.]}
\end{center}
\footnotetext{The corresponding notice in the Academy report of 1849, p. 218, reads: “Mr. Dirichlet read on the determination of mean values in number theory.” K.}

Although the functions considered in number theory are almost never representable by analytic expressions and appear to proceed quite irregularly, a greater regularity nevertheless emerges in their mean values the further one follows their sequence. That is, there are certain simple expressions which represent the progress of the mean values with continuously increasing accuracy, just as one curve approaches the course of another curve of which it is the asymptote. Toward the end of the fifth section of the \emph{Disquisitiones arithmeticae} one finds several highly remarkable expressions of this kind, relating to the theory of quadratic forms. Since these interesting results, which are only communicated there incidentally and without proof, have not yet been proved, and since in general no methods have been known for treating similar questions, I began several years ago to look for suitable means. From that earlier work, however, I gave nothing to the public except a few new results\footnote{Cf. pp. 351 and 372 of volume I of this edition of G. Lejeune Dirichlet's works. K.}, because there appeared to me the prospect of essentially simplifying the treatment of such problems by continued efforts, and in particular of making it independent of the integral calculus. Other investigations then drew me away from this subject for a longer time. After later taking it up again, I became convinced that in many cases one arrives at the asymptotic expression for the mean value by quite elementary considerations founded on a very simple transformation of series. For the moment I restrict myself to a series of problems for which the indicated means alone is sufficient. In a subsequent paper I shall concern myself with more difficult problems whose solution requires combining the transformation just mentioned with other auxiliary means.

\section*{1.}

To put the transformation on which the solution of the problems treated in this paper chiefly rests in its proper light, it seems appropriate to begin at once with one of the simplest questions, to carry its solution as far as can be done without that transformation, and then to complete it with its help.

Let $f(n)$ denote the number of divisors of the integer $n$, and consider the problem of determining the so-called summatory term of this function, that is, the sum
\[
 f(1)+f(2)+\cdots+f(n)=F(n).
\]
If $s$ is an integer $\leq n$, then in as many terms of our sum there will be a unit corresponding to the divisor $s$ as there are multiples of $s$ among the numbers $1,2,\ldots,n$. But the number of these multiples is $\floor{n/s}$, if, as everywhere in what follows, we use square brackets to denote the greatest integer contained in the bracketed value. It follows that
\[
 F(n)=\sum_{1}^{n}\floor{\frac ns},
\]
where the summation sign refers to $s$. This equation immediately gives an approximate determination, that is, an asymptotic expression for $F(n)$. Since $n/s$ exceeds the integer $\floor{n/s}$ by less than one unit, one has, up to an error that cannot exceed $n$,
\[
 F(n)=n\sum_{1}^{n}\frac1s.
\]
But
\[
 \sum_{1}^{n}\frac1s=\log n+C+\frac1{2n}+\cdots,
\]

where $C=0.5772156\ldots$. Since, however, the last expression for $F(n)$ is already affected by an error of order $n$, only the first term of the infinite series is to be retained; and one sees that the equation
\[
 F(n)=n\log n
\]
is exact only up to an error of first order, for which, incidentally, one can easily give a bound that it cannot exceed. Whether the function, which is of order $n\log n$, contains in its asymptotic expression a term of order $n$ with constant coefficient, or, in other words, whether
\[
 \frac1nF(n)-\log n
\]
approaches a fixed limit as $n$ increases, cannot be decided in this way.

\section*{2.}

The transformation already mentioned, with which we now wish to occupy ourselves, rests on the simple observation that, while the terms of the sequence
\[
 \floor{\frac n1},\quad \floor{\frac n2},\quad \ldots,\quad \floor{\frac ns},\quad \ldots,
\]
which stops with the $n$th term, at first decrease very rapidly, from a certain point on each term is either equal to the following one or exceeds it by one unit. This will plainly be the case as soon as the difference
\[
 \frac ns-\frac n{s+1}=\frac{n}{s(s+1)}
\]
has become equal to one or to a proper fraction. Thus if one determines the least integer $\mu$ satisfying
\[
 \mu(\mu+1)\geq n,
\]
the mentioned property will occur at the latest from the $\mu$th term inclusive. Since for our purpose it is quite immaterial whether one more or one fewer term is drawn into the intended transformation, we shall, for greater simplicity, choose for $\mu$ the integer that is either equal to $\sqrt n$ or, when this is irrational, lies immediately above it. Thus
\[
 \mu^2\geq n,
\]
and consequently $\mu(\mu+1)>n$. Put
\[
 \floor{\frac n\mu}=\nu.
\]

Then, as is easy to see, $\nu$ can differ from $\mu$ by at most two units. Let $t$ be any of the numbers $\nu,\nu-1,\ldots,2,1$. It is then easy to determine the index $s$ of the term furthest from the beginning for which $\floor{n/s}$ has the value $t$, and after which a term with value $t-1$ follows. The desired index is given by the double condition
\[
 \frac ns\geq t,
 \qquad \frac n{s+1}<t,
\]
and hence
\[
 s\leq \frac nt,
 \qquad s+1>\frac nt,
 \qquad \text{that is,}\quad s=\floor{\frac nt}.
\]

Now consider the sum
\[
 \floor{\frac n1}\varphi(1)+\floor{\frac n2}\varphi(2)+\cdots+
 \floor{\frac ns}\varphi(s)+\cdots+
 \floor{\frac np}\varphi(p),
\]
where $p$ is, of course, assumed to be greater than $\mu$ and not greater than $n$, or rather consider only that part of the sum which extends from the $\mu$th term onward:
\[
 \floor{\frac n\mu}\varphi(\mu)+\cdots+\floor{\frac ns}\varphi(s)+\cdots+\floor{\frac np}\varphi(p).
\]
We may transform this by combining into partial sums those terms for which $\floor{n/s}$ has the same value, and then adding all the partial sums.

For greater uniformity, omit from the first partial sum the term corresponding to $s=\mu$, join it to the already separated first $\mu-1$ terms, and put
\[
 \floor{\frac np}=q.
\]
Then the partial sums extend respectively from
\[
\begin{array}{llll}
s=\mu & \text{excluded to } s=\floor{n/\nu} & \text{included, with corresponding value } \floor{n/s}=\nu,\\[0.4em]
s=\floor{n/\nu} & \text{to } s=\floor{n/(\nu-1)} && \nu-1,\\[0.4em]
s=\floor{n/(\nu-1)} & \text{to } s=\floor{n/(\nu-2)} && \nu-2,\\
\cdots &&& \cdots,\\[0.4em]
s=\floor{n/(q+1)} & \text{to } s=p && q.
\end{array}
\]

Let $\psi(s)$ denote the summatory term of the function $\varphi(s)$, so that
\[
 \psi(s)=\sum_{h=1}^{h=s}\varphi(h).
\]
The values of the partial sums are then
\[
 \nu\left(\psi\floor{\frac n\nu}-\psi(\mu)\right),\quad
 (\nu-1)\left(\psi\floor{\frac n{\nu-1}}-\psi\floor{\frac n\nu}\right),
\]
\[
 (\nu-2)\left(\psi\floor{\frac n{\nu-2}}-\psi\floor{\frac n{\nu-1}}\right),\quad\ldots,\quad
 q\left(\psi(p)-\psi\floor{\frac n{q+1}}\right),
\]
and their combination gives the expression
\[
 -\nu\psi(\mu)+\psi\floor{\frac n\nu}+
 \psi\floor{\frac n{\nu-1}}+\cdots+
 \psi\floor{\frac n{q+1}}+q\psi(p).
\]
Thus we obtain the general transformation formula
\begin{equation}
 \sum_{1}^{p}\floor{\frac ns}\varphi(s)=
 q\psi(p)-\nu\psi(\mu)+\sum_{1}^{\mu}\floor{\frac ns}\varphi(s)+
 \sum_{q+1}^{\nu}\psi\floor{\frac ns}. \tag{a}
\end{equation}

The case that occurs most frequently in the application of this equation is that in which $p=n$ and consequently $q=1$. Under this assumption one obtains
\begin{equation}
 \sum_{1}^{n}\floor{\frac ns}\varphi(s)=
 -\nu\psi(\mu)+\sum_{1}^{\mu}\floor{\frac ns}\varphi(s)+
 \sum_{2}^{\nu}\psi\floor{\frac ns}. \tag{b}
\end{equation}
As one sees, the advantage of this transformation is that by it a series whose number of terms is $n$, and which contains in each term an expression of the form $\floor{n/s}$, is reduced to two others in which the number of terms still containing $\floor{n/s}$ is lowered to order $\sqrt n$. A similar transformation can still be carried out if, in the expression $\floor{n/s}$, the denominator $s$ is replaced by a function of $s$ increasing with $s$. Since such generality is superfluous for our present purpose, we restrict ourselves to the special form of series just considered.

\section*{3.}

Returning now to the problem treated in Article 1, we obtain, when in equation (b) $\varphi(s)=1$ and hence $\psi(s)=s$, the formula
\[
 F(n)=\sum_{1}^{n}\floor{\frac ns}=-\mu\nu+
 \sum_{1}^{\mu}\floor{\frac ns}+\sum_{1}^{\nu}\floor{\frac ns}.
\]

If in the two sums one puts $n/s$ in place of $\floor{n/s}$, the error is only of order $\sqrt n$; and since likewise
\[
 n\sum_{1}^{\mu}\frac1s=n\sum_{1}^{\nu}\frac1s=\frac12 n\log n+Cn,
 \qquad \mu\nu=n,
\]
it follows, exactly up to an error of order $\sqrt n$, that
\[
 F(n)=n\log n+(2C-1)n.
\]
From the asymptotic expression just found for $F(n)$ one can now easily derive an expression for the mean value of the number of divisors. Write the equation, for greater clarity, in the form
\[
 F(n)=n\log n+(2C-1)n+\xi\sqrt n,
\]
where $\xi$ is unknown but, for however large $n$ may be, remains numerically below a fixed bound that could easily be given. Replace $n$ by $n+s$, with $\xi$ becoming $\xi'$, and divide the difference of the two equations by $s$. One then obtains for the arithmetic mean of the numbers of factors corresponding to the $s$ numbers $n+1,n+2,\ldots,n+s$:
\[
 \log(n+s)+2C-1+\frac ns\log\left(1+\frac sn\right)
 +\frac{\xi'\sqrt{n+s}-\xi\sqrt n}{s}.
\]
If one now lets $s/\sqrt n$ increase beyond every bound, the last term tends to zero; that is, the difference between the mentioned arithmetic mean and the expression formed from the remaining terms becomes smaller than any assignable magnitude. Combining the assumption already made with the further one that $n/s$ also exceeds every bound, one obtains, for the mean corresponding to this double condition, the very simple asymptotic value
\[
 \log n+2C.
\]
This also remains valid when $n$, instead of denoting the number immediately preceding the sequence of $s$ numbers over which the mean is taken, coincides with one of these numbers. For if $m$ is one of them, then $n+t=m$ with $t\leq s$, and under the above assumption the difference $\log m-\log n$ tends to zero.

\section*{4.}

The equation obtained above,
\[
 \sum_{1}^{n}\floor{\frac ns}=n\log n+(2C-1)n,
\]
in which the error is of order $\sqrt n$, gives occasion for an observation at which we shall pause for a moment. Since, on the other hand,
\[
 \sum_{1}^{n}\frac ns=n\log n+Cn
\]
where the error for every $n$ does not exceed a fixed bound, one has, with an error of order $\sqrt n$,
\[
 \sum_{1}^{n}\left(\frac ns-\floor{\frac ns}\right)=(1-C)n.
\]
Thus the arithmetic mean of the values of the difference
\[
 \frac ns-\floor{\frac ns},
\]
corresponding to $s=1,2,\ldots,n$, is equal to $1-C$ and hence is less than $\tfrac12$. It is therefore natural to conjecture that, when $n$ is divided successively by the numbers just named, the case in which the remainder lies below half the divisor will occur more often than the opposite case, in which it is equal to or exceeds half the divisor. We shall test the correctness of this conjecture and try to determine the ratio in which the numbers $1,2,\ldots,n$ divide into two groups in this respect.

The number $s$ gives the first or second case according as
\[
 \frac ns-\floor{\frac ns}<\frac12\quad \text{or}\quad \geq\frac12.
\]
Accordingly one has, respectively,
\[
 \floor{\frac{2n}{s}}-2\floor{\frac ns}=0\quad \text{or}\quad
 \floor{\frac{2n}{s}}-2\floor{\frac ns}=1.
\]
Therefore the number of integers contained in the second group is given by
\[
 \sum_{1}^{n}\floor{\frac{2n}{s}}-2\sum_{1}^{n}\floor{\frac ns},
\]
whose second term has already been determined. For the determination of the first one uses equation (a), in which $n$ is replaced by $2n$ and $p=n$, $q=2$, $\varphi(s)=1$, $\psi(s)=s$ are put. Thus one obtains

\[
 \sum_{1}^{n}\floor{\frac{2n}{s}}=2n-\mu\nu+
 \sum_{1}^{\mu}\floor{\frac{2n}{s}}+
 \sum_{3}^{\nu}\floor{\frac{2n}{s}}.
\]
Since $\mu$ is the integer immediately above $\sqrt{2n}$ and $\nu=\floor{2n/\mu}$, $\mu\nu$ differs from $2n$ only by a quantity of order $\sqrt n$, and the first two terms cancel. The value of the first sum is, still neglecting terms of order $\sqrt n$, given by
\[
 \sum_{1}^{\mu}\floor{\frac{2n}{s}}=2n\sum_{1}^{\mu}\frac1s=2n(\log\mu+C)=n\log 2n+2Cn.
\]
The same value would also hold for the second sum if the two terms corresponding to $s=1$ and $s=2$ were not missing. Therefore, to obtain $\sum_1^n\floor{2n/s}$, one must double the expression just found and subtract
\[
 \floor{\frac{2n}{1}}+\floor{\frac{2n}{2}}=3n.
\]
Thus the number of integers in the second group is found to be
\[
 (\log4-1)n,
\]
and for the first group
\[
 (2-\log4)n.
\]
For large $n$, therefore, the number of divisors having the first property is to the number having the second as $2-\log4$ is to $\log4-1$.

\section*{5.}

Now consider the function $f(n)$ which expresses the sum of the divisors of $n$, or rather, first again as in No. 1, the sum formed from it,
\[
 F(n)=f(1)+f(2)+\cdots+f(n).
\]
By arguments quite similar to those used there, one obtains
\[
 F(n)=\sum_{1}^{n}s\floor{\frac ns}.
\]
Application of equation (b) gives
\[
 F(n)=-\frac12(\mu^2+\mu)\nu+
 \sum_{1}^{\mu}s\floor{\frac ns}+\frac12\sum_{1}^{\nu}\left(\floor{\frac ns}^{2}+\floor{\frac ns}\right).
\]

To see the order of magnitude of the quantities that must be neglected as not completely determinable, first consider the first part of the last sum. Put
\[
 \floor{\frac ns}=\frac ns-\varepsilon,
\]
where $\varepsilon$ is a positive proper fraction depending on $n$ and $s$. Then
\[
 \sum_{1}^{\nu}\floor{\frac ns}^{2}=
 \sum_{1}^{\nu}\frac{n^2}{s^2}-2\sum_{1}^{\nu}\frac ns\varepsilon+
 \sum_{1}^{\nu}\varepsilon^2.
\]
The second and third terms, whose orders cannot exceed $n\log n$ and $\sqrt n$, respectively, must be omitted because of the analytically inexpressible fraction $\varepsilon$ occurring in them. The second part of the sum, namely $\sum_1^{\nu}\floor{n/s}$, which we have already determined and which is of order $n\log n$, is likewise not to be taken into account, although this part can be specified exactly down to and including the first order. Further, neglecting the first order,
\[
 (\mu^2+\mu)\nu=n^{3/2},\qquad
 \sum_{1}^{\mu}s\floor{\frac ns}=n^{3/2};
\]
hence, exactly up to an error of order $n\log n$,
\[
 F(n)=\frac12 n^{3/2}+\frac12 n^2\sum_{1}^{\nu}\frac1{s^2}.
\]
But by a known formula,
\[
 \sum_{1}^{\nu}\frac1{s^2}=\frac{\pi^2}{6}-\frac1\nu+\frac{1}{2\nu^2}-\cdots,
\]
and therefore, excluding the first order,
\[
 \frac12 n^2\sum_{1}^{\nu}\frac1{s^2}=\frac{\pi^2}{12}n^2-\frac12 n^{3/2}.
\]
Consequently one finally obtains
\[
 F(n)=\frac{\pi^2}{12}n^2,
\]
where the error does not exceed order $n\log n$.

Using the expression just found to determine the mean value of $f(n)$, one finds for this mean value, namely for
\[
 \frac1s\left(f(n+1)+\cdots+f(n+s)\right),
\]

the asymptotic expression
\[
 \frac{\pi^2}{6}n,
\]
provided that the simultaneous growth of $s$ and $n$ is understood in such a way that
\[
 \frac ns\quad \text{and}\quad \frac{s}{\log n}
\]
exceed every finite bound. This result, however, has an essentially different meaning from the one found at the end of No. 3. There the difference between the true mean value and its asymptotic expression became smaller than any given magnitude; here this holds only for the ratio of this difference to the true or approximate mean value.

\section*{6.}

In the problems treated so far, to which it seems superfluous to add others solved in quite the same way, the function to be approximately determined had directly the form of the series (a). In other cases this function is given by an equation containing such a series, in whose general term the function to be determined occurs; thus one knows only a recurrence relation between successive values of the function. The simplest example of this kind is given by the function $\varphi(n)$, which expresses the number of integers in the sequence $1,2,\ldots,n$ relatively prime to $n$ and which plays so large a role in number theory. It is well known that this function is expressed by
\[
 \varphi(n)=n\left(1-\frac1a\right)\left(1-\frac1b\right)\left(1-\frac1c\right)\cdots,
\]
where $a,b,c,\ldots$ denote the distinct prime numbers dividing $n$. This formula, however, is not suited to our investigation, and we must choose a different starting point. We begin with the familiar equation
\[
 \sum \varphi(d)=n,
\]
where the summation sign extends over all divisors $d$ of $n$. Adding this equation and the similar ones for $n-1,n-2,\ldots,1$, the term $\varphi(s)$ on the left, where $s\leq n$,

occurs as many times as there are multiples of $s$ in the sequence $1,2,\ldots,n$, that is, $\floor{n/s}$ times. Thus
\[
 \sum_{1}^{n}\floor{\frac ns}\varphi(s)=\frac12 n^2+\frac12 n.
\]

Before going further, let us return for a moment to formula (b) and observe that, for some problems, including the one treated in No. 5, it gives a shortcut when the transformation embraces the whole series, although this sacrifices the advantage, indispensable in other cases, of lowering the number of terms to a smaller order. To obtain the modified formula, observe that, if $t$ denotes quite generally one of the numbers $1,2,3,\ldots,n$, this value is not always found among the terms
\[
 \floor{\frac n1},\quad \floor{\frac n2},\quad \ldots,\quad
 \floor{\frac ns},\quad \ldots,\quad \floor{\frac nn},
\]
since these terms decrease very rapidly at the beginning of the sequence; nevertheless there is always one and only one term that is not greater than $t$ and is followed by another whose value is smaller than $t$. As before, the index $s$ of this term must satisfy
\[
 \frac ns\geq t,
 \qquad \frac n{s+1}<t,
\]
whence
\[
 s=\floor{\frac nt}.
\]
Thus in general $\floor{n/s}=t$ from $s=\floor{n/(t+1)}$ excluded up to $s=\floor{n/t}$ included; and the sum of all terms in which $\floor{n/s}$ has the value $t$ is
\[
 \left(\psi\floor{\frac nt}-\psi\floor{\frac n{t+1}}\right)t.
\]
This expression vanishes and remains correct when no such terms exist, that is, when
\[
 \floor{\frac nt}=\floor{\frac n{t+1}}.
\]
Combining the values of the expression for $t=1,2,\ldots,n$, and observing that for $t=n$ the term $\psi\floor{n/(t+1)}$ is plainly to be reduced to zero, one obtains
\begin{equation}
 \sum_{1}^{n}\floor{\frac ns}\varphi(s)=\sum_{1}^{n}\psi\floor{\frac ns}. \tag{c}
\end{equation}

This transformation is the one we shall now use in our problem. With its help our preceding equation becomes
\[
 \sum_{1}^{n}\psi\floor{\frac ns}=\frac12 n^2+\frac12 n.
\]
It is soon clear that the asymptotic expression for $\psi(n)$ will have the form $\alpha n^2$, where $\alpha$ is a constant. One could at first leave this constant undetermined in the following proof, in which case the value $\alpha=3/\pi^2$ would emerge in the course of the development. Since this value is also easy to foresee, however, we shall for brevity at once consider the expression $\frac3{\pi^2}n^2$. Whatever the still unknown function $\psi(n)$ may be, we may set in general
\[
 \psi(n)=\frac3{\pi^2}n^2+\xi\chi(n),
\]
where $\xi$ also depends on $n$ and the function $\chi(n)$ is arbitrary except for the restriction that it always remain positive, increase with its argument, and not vanish at $n=1$. Suppose $\chi(n)$ has been chosen in the indicated way, and then substitute all values from $n=1$ to $n=N$ into our equation. The corresponding values of $\xi$ will all be finite. Let $A$ be the greatest numerical value so obtained for $\xi$. With this assumed, we bring our equation into the form
\[
 \psi(n)=-\sum_{2}^{n}\psi\floor{\frac ns}+\frac12 n^2+\frac12 n
\]
and consider the values of $n$ lying between $N+1$ and $2N$. Since $\floor{n/s}$ in the sum does not exceed the bound $N$, the part of our sum arising from substituting the second term of the expression assumed above for $\psi(n)$ is numerically smaller than
\[
 A\sum_{2}^{n}\chi\left(\frac ns\right).
\]
The part arising from the first term, namely
\[
 \frac3{\pi^2}\sum_{2}^{n}\floor{\frac ns}^{2},
\]

becomes, if one again writes $\floor{n/s}=n/s-\varepsilon$,
\[
 \frac{3n^2}{\pi^2}\sum_{2}^{n}\frac1{s^2}
 -\frac{6n}{\pi^2}\sum_{2}^{n}\frac{\varepsilon}{s}
 +\frac3{\pi^2}\sum_{2}^{n}\varepsilon^2.
\]
Since
\[
 \sum_{2}^{n}\frac1{s^2}=\frac{\pi^2}{6}-1+\tau,
\]
where $\tau$ is of order $1/n$, and since the two other sums cannot exceed the orders $\log n$ and $n$, respectively, one obtains an equation
\[
 \psi(n)=\frac3{\pi^2}n^2+\xi,
\]
in which $\xi$ is numerically smaller than
\[
 Pn\log n+A\sum_{2}^{n}\chi\left(\frac ns\right),
\]
where $P$ is a sufficiently large constant independent of $N$. Comparing this result, valid from $n=N+1$ to $n=2N$, with the expression assumed above for $\psi(n)$,
\[
 \frac3{\pi^2}n^2+\xi\chi(n),
\]
shows that for this new interval the greatest numerical value of $\xi$ does not exceed the maximum, in that interval, of the expression
\[
 \frac{Pn\log n}{\chi(n)}+\frac{A}{\chi(n)}\sum_{2}^{n}\chi\left(\frac ns\right).
\]
If one now gives the previously undetermined function $\chi(n)$ the form of a positive power $n^\delta$, then the expression multiplied by $A$ becomes
\[
 \sum_{2}^{n}\frac1{s^\delta}<\sum_{2}^{\infty}\frac1{s^\delta}=q,
\]
and $\delta$ can be chosen between 1 and 2 so that the constant $q$ is a proper fraction. On the other hand, $N$ can be chosen so large that for $n\geq N$,
\[
 \frac{Pn\log n}{n^\delta}<k,
\]
where the constant $k$ is arbitrarily small. If $A'$ denotes the greatest numerical value of $\xi$ occurring between $n=N+1$ and $n=2N$, then
\[
 A'<Aq+k.
\]

Since $k$ can be made arbitrarily small as $N$ increases, while $A$ does not decrease and $q$ remains constant, for sufficiently large $N$ one has
\[
 A'<A.
\]
That is, the maximum $A$ originally valid up to $n=N$ is also valid up to $2N$, and for the same reason also up to $4N,8N,\ldots$, and therefore quite generally.

Thus
\[
 \psi(n)=\frac3{\pi^2}n^2
\]
with an error that cannot exceed order $n^\delta$, where the constant $\delta$ exceeds, however slightly, the value $\gamma$ given by the equation
\[
 \sum_{2}^{\infty}\frac1{s^\gamma}=1.
\]
For the mean value of $\varphi(n)$ this gives the expression
\[
 \frac6{\pi^2}n,
\]
whose meaning is clear from the foregoing.

\section*{7.}

As a final example we choose the function $\varphi(n)$ which denotes the number of all possible decompositions of $n$ into two factors without a common divisor, and which, as is well known, has the expression $2^\rho$, where $\rho$ is the number of distinct prime divisors of $n$. Put again
\[
 \sum_{1}^{n}\varphi(s)=\psi(n).
\]
Then $\psi(n)$ stands in a simple relation to the function $F(n)$ treated in Articles 1 and 3. Indeed, $F(n)$ plainly expresses the number of pairs of integers $x,y$ satisfying $xy\leq n$, whereas $\psi(n)$ denotes the number of pairs of relatively prime integers $\xi,\eta$ for which likewise $\xi\eta\leq n$. Divide the pairs $x,y$ into groups, the $s$th containing those pairs for which $s$ is the greatest common divisor; thus, if $x=\xi s$, $y=\eta s$, the numbers $\xi$ and $\eta$ are relatively prime. Dividing the inequality by $s^2$ gives
\[
 \xi\eta\leq\frac n{s^2}\quad \text{or}\quad
 \xi\eta\leq\floor{\frac n{s^2}}.
\]

Conversely, every pair of relatively prime integers $\xi,\eta$ satisfying this condition gives, by multiplication by $s$, two integers $x,y$ with greatest common divisor $s$ and with $xy\leq n$. Hence the number of pairs contained in the $s$th group is expressed by $\psi\floor{n/s^2}$. Thus one obtains the equation
\[
 \sum \psi\floor{\frac n{s^2}}=F(n)=n\log n+(2C-1)n,
\]
where the sum extends from $s=1$ to $s=\floor{\sqrt n}$, and by the above the neglected term does not exceed order $\sqrt n$. It is consequently easy to see that the asymptotic expression for $\psi(n)$ will have the form
\[
 \alpha n\log n+\beta n,
\]
where $\alpha$ and $\beta$ are two constants still to be determined. If, as in the preceding article, one sets in our equation
\[
 \psi(n)=\alpha n\log n+\beta n+\xi n^\delta,
\]
and chooses $\alpha$ and $\beta$ so that the terms of orders $n\log n$ and $n$ cancel, namely
\[
 \alpha=\frac6{\pi^2},
 \qquad
 \beta=\frac6{\pi^2}\left(\frac{12C'}{\pi^2}+2C-1\right),
\]
where $C$ has its former meaning and
\[
 C'=\sum_{2}^{\infty}\frac{\log s}{s^2},
\]
then by arguments entirely analogous to those developed in the preceding article one finds that $\xi$, however large $n$ may be, remains under a fixed bound, provided only that $\delta>\frac12\gamma$, where $\gamma$ is understood in the above sense. From the expression thus obtained for $\psi(n)$,
\[
 \alpha n\log n+\beta n,
\]
there follows for $\varphi(n)$, in the sense determined by what has just been said, the mean value
\[
 \frac6{\pi^2}\left(\log n+\frac{12C'}{\pi^2}+2C\right).
\]
The question just treated is connected with the mean value given in Article 301 of the \emph{Disquisitiones arithmeticae} for the number of \emph{genera} corresponding to a negative determinant $-n$. According to known theorems, the expression for the number of \emph{genera}, corresponding to the five linear forms
\[
 n=8h,\quad n=8h+4,\quad n=4h+2,\quad n=4h+1,\quad n=4h+3,
\]

is respectively
\[
 \varphi(n),\quad \tfrac12\varphi(n),\quad \tfrac12\varphi(n),\quad
 \varphi(n),\quad \tfrac12\varphi(n).
\]
On the other hand, by obvious modifications of the preceding considerations, whose execution we leave to the reader, one can determine the approximate value of $\psi(n)$ when $n$ has a prescribed linear form, and the sum $\sum\varphi(s)=\psi(n)$ is no longer extended over all the numbers $1,2,\ldots,n$, but only over the numbers of this form contained in that sequence; from this one can then derive the mean value of $\varphi(n)$. For the linear forms listed above,
\[
 n=8h,\quad n=8h+4,\quad n=4h+2,\quad n=4h+1,\quad n=4h+3,
\]
one finds as the mean value of $\varphi(n)$, if for abbreviation one puts
\[
 \log n+\frac{12C'}{\pi^2}+2C=\Delta,
\]
the following five expressions:
\[
 \frac8{\pi^2}\left(\Delta-\frac83\log2\right),\quad
 \frac8{\pi^2}\left(\Delta-\frac23\log2\right),\quad
 \frac8{\pi^2}\left(\Delta+\frac13\log2\right),
\]
\[
 \frac4{\pi^2}\left(\Delta+\frac43\log2\right),\quad
 \frac4{\pi^2}\left(\Delta+\frac43\log2\right).
\]
By the foregoing, these expressions, multiplied respectively by $1,\tfrac12,\tfrac12,1,\tfrac12$, give the mean numbers of the \emph{genera} of determinant $-n$ for the previously enumerated linear forms. Since among every eight consecutive numbers, respectively $1,1,2,2,2$ are contained in these forms, the sum of our expressions multiplied successively by
\[
 1\cdot\tfrac18,\quad \tfrac12\cdot\tfrac18,\quad
 \tfrac12\cdot\tfrac28,\quad 1\cdot\tfrac28,\quad
 \tfrac12\cdot\tfrac28
\]
is, that is,
\[
 \frac4{\pi^2}\left(\Delta-\frac12\log2\right),
\]
the desired mean number of \emph{genera} corresponding to the determinant $-n$, when this determinant is thought of generally, that is, no longer restricted to a particular linear form. This agrees with the result given at the cited place. That the same expression also holds for the positive determinant $n$ follows easily from the fact that the mean values of $\varphi(n)$ corresponding to the linear forms $4h+1$ and $4h+3$ coincide, and that, on the other hand, the exceptions occurring for square determinants are plainly without influence on the final result.

\end{document}
